\documentclass[12pt]{article}
\usepackage[T1]{fontenc}
\usepackage[utf8]{inputenc}
\usepackage{lmodern}
\usepackage{textcomp}
\usepackage{enumitem}
\usepackage{geometry}
\usepackage{graphicx}
\usepackage{amsmath, amssymb}
\geometry{margin=1in}
\begin{document}
\begin{center}
Philosophy - Undergraduate Level Assessment 6 \
Total Marks: 40 \
Answer all questions.
\end{center}

\section*{Multiple Choice (10 marks)}
\begin{enumerate}[label=\arabic*.]
\item Berkeley insists that heat and cold are \_\_\_\_\_.
\begin{enumerate}[label=(\alph*)]
    \item elements of nature that do not exist independently
    \item only things existing apart from our minds
    \item only sensations existing in our minds
    \item physical objects
    \item manifestations of our subconscious
\end{enumerate}
\item One suggestion that Lukianoff and Haidt make to challenge vindictive protectiveness is
\begin{enumerate}[label=(\alph*)]
    \item to implement stricter guidelines for classroom discussions.
    \item to promote greater understanding of historical and contemporary oppression.
    \item to increase the number of safe spaces on campus.
    \item to mandate sensitivity training for all students.
    \item to increase funding for mental health services on campus.
\end{enumerate}
\item According to Lukianoff and Haidt, the recent trend to uncover microaggressions encourages
\begin{enumerate}[label=(\alph*)]
    \item students to confabulate reasons.
    \item the pursuit of justice by marking out racism, sexism, and classism.
    \item labeling, by assigning global negative traits to persons.
    \item universities to bear overly burdensome legal obligations.
\end{enumerate}
\item In a hypothetical syllogism, when the minor premise affirms the antecedent
\begin{enumerate}[label=(\alph*)]
    \item no valid conclusion can be drawn
    \item the conclusion must affirm the consequent
    \item the conclusion must deny the consequent
    \item the conclusion must deny the antecedent
\end{enumerate}
\item According to Mill, the value of a particular pleasure depends on
\begin{enumerate}[label=(\alph*)]
    \item its rarity or frequency.
    \item societal norms and values.
    \item the individual's personal preference.
    \item the amount of effort required to obtain it.
    \item the potential pain that might accompany it.
\end{enumerate}
\end{enumerate}

\section*{True / False (10 marks)}
\textit{Answer True or False}
\begin{enumerate}[label=\arabic*.]
\setcounter{enumi}{5}
\item True or False: According to Butler, justice is the only right that individuals possess in a strict philosophical sense.
\item True or False: Singer's principle of equality asserts that all living beings are inherently equal in their moral status and treatment.
\item True or False: Ashford's article addresses the daunting task of solving worldwide economic imbalance.
\item True or False: In Sikh traditions, the phrase Guru-Panth refers to the concept of community among followers and the collective path they share.
\item True or False: In Hobbes's view, to say something is good is to say that it aligns with natural law.
\end{enumerate}

\section*{Long Form Response (20 marks)}
\textit{Answer each question with a clear and complete explanation.}
\begin{enumerate}[label=\arabic*.]
\setcounter{enumi}{10}
\item Discuss Hursthouse's perspective on the controversy surrounding the claim that an adequate moral theory must provide clear, actionable guidance for intelligent individuals.
\item Discuss Pogge's argument regarding the moral responsibilities of individuals in wealthy countries towards the global poor, particularly focusing on the concept of libertarian duties.
\end{enumerate}

\vfill
\noindent\textit{End of Paper}
\end{document}